\documentclass[11pt,a4paper]{article}
\usepackage[margin=25mm]{geometry}
\usepackage{amsmath,amssymb,amsthm,booktabs}
\newtheorem{lemma}{Lemma}\newtheorem{theorem}{Theorem}
\title{The two remaining order 15 counts}
\author{Human Compression project}\date{}
\begin{document}\maketitle
For a finite group $J$, let $x_J(p,q)$ count ordered pairs $(r,s)$
of orders $p,q$ whose product has order two, and let $e_J(p,q)$
count those pairs generating $J$. Put $G=O'N$. The accepted
order 20 argument proves that every cyclic subgroup of order five
has a unique row 11 maximal overgroup, denoted $B$. This follows
from the complete roster and the direct incidence
$17778376\cdot36=640021536$. Thus every pair with an element of
order 15 belongs to at most one row 11 maximal subgroup.
We use that theorem to recover the two remaining entries in the
order 15 column of the fixed Connor--Leemans table.

\begin{lemma}\label{projective}
Every pair counted by $x_{\operatorname{PSL}_2(31)}(4,15)$
generates $\operatorname{PSL}_2(31)$.
\end{lemma}
\begin{proof}
For each of the two global classes, rows 15 and 16, the complete
maximal list consists of a subgroup of order 465, two classes
of $A_5$, a dihedral group of order 32, a dihedral group of
order 30, and two classes of order 24. Their orders exclude
the required pair except for $A_5$. An even permutation of five
letters has no order 15: disjoint cycles of lengths three and
five would require eight letters. Hence no proper subgroup
contains such a pair.
\end{proof}

\begin{lemma}\label{janko}
For $J=J_1$,
\[
e_J(6,15)=x_J(6,15)-2926\cdot40.
\]
Every subtracted pair lies in a row 11 maximal subgroup of $G$.
\end{lemma}
\begin{proof}
The complete maximal list of $J$ has orders
$660,168,120,114,110,60,42$. The order 660 subgroup is
$\operatorname{PSL}_2(11)$, whose nonidentity element orders
divide five or six, or equal eleven; it has no element of
order 15. The order 120 subgroup is $A_5\times C_2$, also
without such an element. The remaining orders leave only
row 366, $E=D_{10}\times D_6$, of which there are 2926 in $J$.
Here $D_{2n}$ denotes a dihedral group of order $2n$.

Write $D_{10}=\langle u,w\mid u^5=w^2=1,\ wuw=u^{-1}\rangle$
and $D_6=\langle v,z\mid v^3=z^2=1,\ zvz=v^{-1}\rangle$.
The order 15 elements of $E$ are exactly
$s=(u^i,v^j)$ with $i\in\{1,2,3,4\}$ and $j\in\{1,2\}$.
An element of order six has the form $r=(u^aw,v^k)$, with
$a\in\mathbb Z/5\mathbb Z$ and $k\in\{1,2\}$.
The product has order two exactly when $k=-j$ modulo three.
There are consequently $8\cdot5=40$ pairs.
Each generates the same subgroup
\[
H=D_{10}\times\langle v\rangle,
\]
of order 30. Indeed $r^3$ and $s^3$ generate the first dihedral
factor, and the second factor then follows from $r$.
This subgroup is normal in $E$.

The roster has exactly three classes of order 30, rows 463--465.
Only row 464 contains both a cyclic subgroup of order six and
one of order fifteen. The class of $H$ is therefore determined
without interpreting an ambiguous structure label. Its class
length is 2560086144, so $|N_G(H)|=180$. Since $|J|=175560$,
\[
|N_J(H)|\mid\gcd(180,175560)=60.
\]
On the other hand $E\leq N_J(H)$ has order 60. Hence
$N_J(H)=E$. Two distinct row 366 subgroups of this fixed $J$
cannot share their generated subgroup $H$. All 2926 collections
of 40 pairs are therefore disjoint, proving the formula.

Finally $E$ normalizes its cyclic rotation subgroup $\langle u\rangle$
of order five. By uniqueness of its row 11 maximal overgroup,
$E$ normalizes that maximal subgroup. A proper maximal subgroup
of the simple group $G$ is self-normalizing. Thus $E$ lies in
that row 11 maximal subgroup, as claimed.
\end{proof}

\begin{theorem}
The ordered automorphism-orbit counts of generating pairs are
\[
\frac{e_G(4,15)}{|\operatorname{Aut}(G)|}=234,
\qquad
\frac{e_G(6,15)}{|\operatorname{Aut}(G)|}=874.
\]
\end{theorem}
\begin{proof}
The acquired complete maximal character tables have nonzero
$(4,15)$ contributions only in rows 11, 15 and 16. The last
two consist entirely of pairs generating their maximal subgroup,
by Lemma~\ref{projective}, so cannot overlap another maximal
contribution. The row 11 contribution is unique by the initial
order-five argument. It follows that
\[
e_G(4,15)=x_G(4,15)-\sum_M[G:M]x_M(4,15).
\]
For $(6,15)$ only rows 6 and 11 contribute; row 6 is $J_1$.
The proper pairs in $J_1$ are all already covered by row 11,
whereas a pair generating $J_1$ belongs to no other maximal.
Lemma~\ref{janko} therefore gives
\[
e_G(6,15)=x_G(6,15)-\sum_M[G:M]x_M(6,15)
              +2624832\cdot2926\cdot40.
\]
Substituting the acquired exact character evaluations and dividing
by $|\operatorname{Aut}(G)|=2|G|$ gives the stated values. The
action on generating pairs is free; no duality quotient is taken.
Both answers agree with the fixed printed entries.
\end{proof}

\paragraph{Attribution and verification.}
The numerical assertions are from Connor and Leemans,
\emph{Algorithmic enumeration of regular maps}, Table~1.
The complete roster, its subgroup identifications and the acquired
character tables are borrowed premises. We reuse the accepted
unique order-five overgroup theorem. The present argument proves
that the normalizer inside a fixed $J_1$ is exactly the dihedral
product, despite the larger normalizer in $G$. A literal reader
also checks all 60 elements of the explicit product, all 40
pairs and their complete generated closures. The two computations
here give numerical counts only; they do not settle remaining
chirality assertions or complete the original manuscript.
\end{document}
